\chapter{数学与计算机：算法能走多远}

\section{一千万次成功之后}

先来做一个小游戏。随便写一个正整数，然后反复执行下面的规则：如果这个数是偶数，就把它除以 $2$；如果是奇数，就把它乘 $3$ 再加 $1$。比如从 $6$ 出发：

\[6 \to 3 \to 10 \to 5 \to 16 \to 8 \to 4 \to 2 \to 1.\]

到了 $1$ 之后，$1$ 会变成 $4$，$4$ 变回 $2$，$2$ 又变回 $1$，陷入 $4 \to 2 \to 1$ 的小圈子，我们就停在这里。

这个规则简单到不能再简单，可它提出的问题却难倒了全世界。数学家们猜想：\textbf{无论}从哪个正整数出发，经过有限步之后都会掉进 $4 \to 2 \to 1$ 的圈子。这就是著名的"$3n+1$ 猜想"，也叫角谷猜想或冰雹猜想（因为数在降下来之前常常忽上忽下，像冰雹在云里翻滚）。

这个猜想至今没有被证明，也没有被推翻。但是人们做了一件很"笨"的事：让计算机去试。一台普通的笔记本电脑，一秒钟就能检查几百万个数；专业团队用上大型计算机，已经把从 $1$ 到 $2^{68}$ 左右——大约三亿亿——之间的所有整数都验证了一遍。\textbf{每一个数}最终都落进了 $4 \to 2 \to 1$ 的圈子。一个例外都没有。

三亿亿次成功，无一例外。现在问你一个问题：这个猜想算是被证明了吗？

直觉上，很多同学会说"肯定是对的呀，都试了这么多了"。也有人会犹豫："万一在更大更大的地方藏着一个反例呢？"这一章，我们就从这个犹豫出发，去看看数学和计算机之间这段既亲密又微妙的关系：计算机到底能为数学做什么，又永远做不到什么。

\section{验证了这么多，为什么还不算证明？}

先别急着站队，我们来看几个历史故事，看看"试了很多次都对"这件事有多么靠不住。

\subsection{一个骗了四十次的公式}

考虑式子 $n^2 - n + 41$。把 $n=1$ 代进去，得 $41$，是素数；$n=2$ 得 $43$，素数；$n=3$ 得 $47$，还是素数。你可以继续代：$n=4,5,6,\dots$ 一路代下去，得到的 $43,47,53,61,71,83,97,\dots$ 竟然全都是素数！事实上，从 $n=1$ 一直代到 $n=40$，整整四十次，次次都是素数。

如果让计算机来验证，它会信心满满地报告："连续四十次成功！" 可是 $n=41$ 的时候：

\[41^2 - 41 + 41 = 41^2 = 41 \times 41,\]

这显然不是素数。四十次成功之后紧跟着的就是失败。

\subsection{藏在九亿后面的反例}

如果说四十次太少，那九亿次够不够？数学家波利亚在 1919 年提出过一个关于素数分布的猜想（细节不重要，重要的是它是一个"看起来大概率正确"的命题）。人们逐个验证，发现它对九亿以内的所有数都成立。九亿个数啊，一个人一秒钟检查一个，要不吃不喝查上二十多年。

然而 1958 年有人证明，这个猜想是错的，最小的反例出现在 $906\,150\,257$ 附近——九亿多一点的地方。换句话说，前九亿次验证全部"成功"，却丝毫没有触碰到真相的边缘。大自然把第一个例外藏在了一个靠蛮力极难够到的地方。

\subsection{大人物的猜想也会翻船}

大数学家欧拉在 1769 年提出过一个猜想，大意是：至少要有 $n$ 个"$n$ 次方的数"加起来，才可能等于另一个数的 $n$ 次方。这个猜想挺立了近两百年，很多人都相信它——毕竟欧拉看问题极少出错。

结果到了 1966 年，两位数学家借助计算机搜索，找到了一个干干净净的反例：

\[27^5 + 84^5 + 110^5 + 133^5 = 144^5.\]

按照欧拉的猜想，至少需要 $5$ 个五次方才能凑出一个五次方，可这里只用了 $4$ 个。近两百年的信任，被一行等式推翻。值得注意的是，是\textbf{计算机}帮助找到了这个反例——机器搜索了人工不可能搜索的庞大范围。计算机在这里扮演的角色不是"证明者"，而是"找茬高手"。

\begin{fanli}
\textbf{反例警报：验证了 $N$ 次，就宣布结论成立。}

不管 $N$ 是一百、一百万还是三亿亿，"前 $N$ 个都对"只说明结论在\textbf{有限范围}内成立。数学猜想谈论的通常是\textbf{无穷多个}对象，而无穷是试不完的。有限次验证可以提供\textbf{证据}，让我们更有信心；可以帮我们\textbf{发现规律}；运气好的话还能\textbf{撞上反例}把猜想推翻——但它永远不能代替覆盖所有情况的证明。反过来说，推翻一个猜想只需要一个反例，而证明一个猜想需要对付无穷多种情况，这两件事的难度天生就不对称。
\end{fanli}

回到 $3n+1$ 猜想：三亿亿次验证给了我们很强的信心，但波利亚猜想的故事提醒我们，第一个例外完全可能藏在一百亿亿亿之后。所以在它被真正证明（或者找到反例）之前，我们只能老老实实地叫它"\textbf{猜想}"，而不是定理。

\section{算法：一台思想上的机器}

刚才我们一直在说"让计算机去验证"。可是计算机到底在干什么？它既没有眼睛也没有直觉，它做的事本质上只有一件：\textbf{一步一步地执行预先写好的指令}。这样的指令序列有个名字——算法。

\begin{shuyu}{算法}
\textbf{直观解释}：算法就是一份写得足够死的操作手册，死到让一个完全不懂数学、只会机械照办的人（或机器）也能一步步执行下去，并且在有限步之后给出结果。

\textbf{正式定义}：算法是一组有限的、明确的指令，对每类合法的输入，按指令执行有限步之后停止，并输出一个确定的答案。"有限步内停止"这一条很关键——永远转个不停的程序不配叫算法。

\textbf{一个例子}：求两个正整数的最大公约数，可以用两千年前的辗转相除法。比如求 $1071$ 和 $462$ 的最大公约数：$1071 = 2 \times 462 + 147$，$462 = 3 \times 147 + 21$，$147 = 7 \times 21 + 0$，余数为 $0$ 时停下，答案是 $21$。每一步都有明确指令，有限步内必停，这就是教科书级别的算法。
\end{shuyu}

为什么我们这么在意"算法"这个概念？因为它给"机械地计算"划了一条清晰的边界。一旦把"算法"定义清楚了，我们就可以问一些惊人的问题：

\begin{itemize}
\item 是不是\textbf{所有}数学问题都有算法能解决？
\item 如果有，我们能不能造一台机器，把任何数学命题输进去，它就自动告诉我们"真"还是"假"？
\end{itemize}

第二个问题听起来简直是数学家的终极梦想——如果这种机器存在，证明定理就成了按按钮的体力活，数学家们可以集体退休了。二十世纪初，确实有不少数学家怀抱着这样的希望。

\section{有些事，算法永远做不到}

梦想在 1936 年被击碎了。英国数学家图灵在思考"计算"的本质时，提出了一个朴素的问题：\textbf{能不能造一个算法，判断任意给定的程序会不会停机？}

先解释一下什么叫"停机"。程序有两种命运：要么运行有限步后结束（停机），比如辗转相除法；要么陷入死循环，永远跑下去，比如一个写着"从 $1$ 开始不断加 $1$，永不停止"的程序。一个程序会不会停机，取决于它的代码和输入。有些程序一眼就能看出会死循环，有些程序——比如"从某个大数开始不断执行 $3n+1$ 规则"——会不会停，恰恰是没人知道的难题。

\begin{dingyi}
\textbf{停机问题}：是否存在一个算法 $H$，给它任意程序 $P$ 和任意输入 $x$，它都能在有限时间内正确回答"$P$ 作用在 $x$ 上会不会停机"？
\end{dingyi}

图灵的答案石破天惊：这样的算法\textbf{不存在}。不是"暂时没人找到"，而是逻辑上根本不可能存在，过去、现在、未来都不会有。这是本章的核心定理。

\begin{dingli}[停机问题不可判定]
不存在一个能对任意程序和任意输入判定其是否停机的通用算法。
\end{dingli}

\begin{proof}
用反证法。假设这样的算法 $H$ 存在。$H$ 接收一段程序文本和一份输入，总是能在有限步后回答"停"或"不停"，而且永远答对。

现在我们用 $H$ 做零件，拼装一个怪程序 $D$。$D$ 只吃一种东西：一段程序文本 $Q$。$D$ 的工作流程是：

\begin{enumerate}
\item 先调用 $H$，问它："程序 $Q$ 拿它\textbf{自己的文本}当输入时，会停机吗？"
\item 如果 $H$ 回答"会停"，$D$ 就故意进入死循环，永不停机；
\item 如果 $H$ 回答"不会停"，$D$ 就立刻停机。
\end{enumerate}

简单说，$D$ 的脾气是：\textbf{你说我停，我偏不停；你说我不停，我偏停}。构造 $D$ 没有任何技术障碍——既然 $H$ 是个现成的算法，把它包进 $D$ 里只是写几行代码的事。

好戏来了：把 $D$ 自己的文本喂给 $D$，问：$D$ 作用在 $D$ 上，停不停机？

\begin{itemize}
\item 假如 $D(D)$ 停机。那 $H$ 对"$D$ 作用在自己上会不会停"的回答必然是"会停"，可按照 $D$ 的设计，$H$ 答"会停"时 $D$ 偏偏要进死循环——$D(D)$ 不停机。矛盾。
\item 假如 $D(D)$ 不停机。那 $H$ 的回答必然是"不会停"，可按照 $D$ 的设计，$H$ 答"不会停"时 $D$ 立刻停机——$D(D)$ 停机。又矛盾。
\end{itemize}

两种情况都导致矛盾，问题只能出在最初的假设上：算法 $H$ 根本不存在。
\end{proof}

这个证明的巧妙之处在于"自己指自己"：让一个判断程序去判断一个专门跟它对着干的程序，于是自相矛盾。这和说谎者悖论（"这句话是假的"）、和前面章节里见过的对角线方法，在精神上是亲戚——\textbf{任何强大到能谈论自身的系统，都会被"自我指涉"咬到尾巴}。

停机问题不可判定，带来的震动远不止"程序员没法写出万能死循环检测器"这么简单。它意味着：\textbf{存在这样的数学问题，答案明明客观存在（程序要么停要么不停，没有第三种可能），却没有任何算法能系统地求出它。}这类问题被称为"不可计算的"或"不可判定的"。后来人们还证明，许多看上去纯粹的数学问题——比如"任给一个多项式方程组，判断它有没有整数解"（希尔伯特第十问题）——同样不可判定。

\begin{lianjie}
\textbf{连接线：停机问题与哥德尔不完备定理。}

在稍早的章节里我们见过哥德尔的不完备定理：任何足够强的数学公理系统里，都存在"真但无法在该系统内证明"的命题。停机问题和它是同一场大地震的两个震感。事实上，如果一个数学系统完备到能判定一切命题，我们马上就能用它造出停机判定算法——把"程序 $P$ 会停机"这个命题交给系统去判就行了。既然停机问题不可判定，完备的系统就不可能存在。计算和证明，在深渊底部是连通的。
\end{lianjie}

\section{快与慢的鸿沟：计算复杂度}

停机问题告诉我们：有些问题\textbf{原则上}就没有算法。接下来一个更实际的问题是：在那些\textbf{有}算法的问题里，是不是所有算法都一样好用？

答案显然是否定的。看两个算法：

\textbf{顺序查找}：在一摞 $n$ 张打乱的书签里找夹着枫叶的那张，只能一张一张翻，最坏情况要翻 $n$ 次。

\textbf{二分查找}：如果书签按页码排好了序，先看正中间那张，页码偏大就往左半摞找，偏小就往右半摞找，每翻一次就排除一半。$n$ 张书签最多翻 $\log_2 n$ 次——$1000$ 张只要 $10$ 次，一百万张也只要 $20$ 次。

同样一个问题，算法不同，工作量天差地别。计算机科学干脆把问题本身按"最好的算法有多快"分类，这就是\textbf{计算复杂度}研究的事。其中最引人注目的一条鸿沟，横在两类问题之间：

\begin{itemize}
\item \textbf{"容易算"的问题}（记作 P 类）：存在算法，所需步数大约是输入规模的\textbf{多项式}级别，比如 $n$、$n^2$、$n^3$。排序、求最大公约数、找最短路径都在这一类。
\item \textbf{"容易验"的问题}（记作 NP 类）：答案可能很难\textbf{找}，但一旦有人把答案送到你面前，你能在多项式时间内\textbf{验证}它对不对。数独是典型：填满一个九宫格可能让你抓耳挠腮，但检查一份填好的答案合不合法，只需扫一遍每行每列每宫。
\end{itemize}

每个 P 类问题天然也是 NP 类（能算出来当然能验证）。反过来呢？\textbf{每个容易验证的问题，是不是其实也容易算？}这就是价值百万美元的"P 对 NP"问题，至今无人解决。大多数专家猜测 P 不等于 NP——也就是说，确实存在"验证容易、寻找难"的问题。

为什么多项式和别的增长差别这么大？看一张表就明白了。假设计算机每秒执行十亿步：

\begin{center}
\begin{tabular}{cccc}
\toprule
输入规模 $n$ & $n^2$ 步 & $2^n$ 步 & $2^n$ 步大约耗时 \\
\midrule
$10$ & $100$ & $1024$ & 一瞬间 \\
$20$ & $400$ & 约一百万 & 一瞬间 \\
$30$ & $900$ & 约十亿 & 约一秒 \\
$40$ & $1600$ & 约一万亿 & 约十八分钟 \\
$60$ & $3600$ & 约一百亿亿 & 约三十六年 \\
\bottomrule
\end{tabular}
\end{center}

多项式（比如 $n^2$）长得不算快，输入翻几倍，时间翻几十倍，咬咬牙总能算完。指数（比如 $2^n$）却是头怪兽：输入每加 $1$，工作量直接翻倍。很多重要的问题——比如"给几十个城市排一条最短的巡回路线"——目前只找得到指数级别的算法。输入稍大一点，宇宙的年龄都不够用。注意这和停机问题不同：这些问题\textbf{有}算法，只是算法慢得没有实用价值。数学把它们和"原则上无解"的问题区分开来，一层一层地给"难"分级。

\begin{center}
\begin{tikzpicture}[scale=0.55]
\draw[->] (0,0) -- (7,0) node[right] {$n$};
\draw[->] (0,0) -- (0,6.5) node[above] {步数};
\draw[thick] (0,0) -- (1,0.25) -- (2,1) -- (3,2.25) -- (4,4) -- (4.9,6.2);
\node at (5.6,4.6) {$n^2$};
\draw[thick,dashed] (0,0.2) -- (1,0.4) -- (2,0.8) -- (3,1.6) -- (4,3.2) -- (4.6,6.4);
\node at (3.2,3.4) {$2^n$};
\foreach \x in {1,...,6} \draw (\x,0.05) -- (\x,-0.05) node[below] {\small \x};
\end{tikzpicture}
\end{center}

\begin{zhu}
上图只是示意草图：两条曲线都在加速，但虚线的指数增长很快就会把任何多项式甩得无影无踪——虽然在小 $n$ 处它们看上去差别不大。这正是"渐近"（看长远趋势）思想的价值。
\end{zhu}

这条鸿沟不只是学术游戏。你每天上网购物、登录账号时用的密码，其安全性就建立在"验证容易、逆向难"之上：把两个大素数乘起来很容易，把一个几百位的大数拆回两个素数却难到所有超级计算机联手都算不动。如果有人证明了 P 等于 NP 并且构造出快速算法，整个现代密码体系的地基就会动摇。

\section{计算机反过来帮助数学}

前面我们强调了计算机的局限：不能靠蛮力验证代替证明，有些问题永远算不出来，有些问题算出来也等不起。但如果就此得出"计算机对数学没什么用"的结论，那就大错特错了。真实的图景有趣得多：计算机正在以至少三种方式深刻地改变数学本身。

\subsection{海量计算，作为证明的一部分：四色定理}

1852 年有人提出：给地图上色，要求相邻国家颜色不同，四种颜色够不够？这个"四色问题"折腾了数学家一百二十多年。1976 年，两位数学家宣布证明了它——但证明的方式前所未闻：他们把问题归结为\textbf{一千多种}基本情形，然后用计算机逐一检查了这一千多种情形全部成立。定理成立，但其中一千多步推理是机器干的，没有任何人能手算复查。

数学界炸开了锅：这算证明吗？没人能人工核对的证明，你敢信吗？争论持续了很多年。如今主流的观点是：只要程序本身被严格审查、计算逻辑清晰，机器完成的穷举可以算证明的一部分。后来这个证明被不断简化、重新验证。但四色定理永远改变了"证明"这个词的含义——\textbf{证明的可靠性，从"每个聪明人都能复查"变成了"逻辑链条的每一环都可被审查，哪怕其中一环只有机器干得动"}。

\subsection{机器当校对员：形式化证明}

另一条路更彻底：与其争论机器参与的证明算不算数，不如干脆把证明\textbf{整个}写成机器能检查的严格语言。这就是"形式化证明"。像 Lean、Coq 这样的"证明助手"软件，要求数学家把每一步推理写到毫无含糊的程度，然后由机器逐行检查逻辑是否严密。人负责出主意，机器负责挑毛病。

这项工作枯燥但回报惊人。最著名的战例是开普勒猜想：把同样大小的球堆进箱子，最密的堆法就是水果店那种堆法。这个命题直觉上"显然"，证明却难如登天。1998 年，数学家黑尔斯宣布证出，证明里有大量计算机计算；审稿人花了几年的工夫，最后只能表示"百分之九十九确信它是对的"——那剩下的百分之一，正是没人查得动的机器计算部分。黑尔斯干脆发起了一个工程，把整个证明用形式化语言重写，2014 年机器盖章确认：每一步逻辑无误。怀疑彻底消失。

\begin{pandeng}
\textbf{继续攀登}

想往高处看，可以留意这些关键词：\textbf{图灵机}（图灵为"算法"建立的数学模型，整个计算机科学的地基）；\textbf{P vs NP 问题}（克雷数学研究所七大千禧年难题之一，悬赏一百万美元）；\textbf{希尔伯特第十问题}（丢番图方程可解性的不可判定性）；\textbf{证明助手 Lean}（如今数学界正在用它把现代数学的整个大厦逐步"机器化"）；\textbf{SAT 求解器}（工业界验证芯片逻辑的超级穷举引擎）。
\end{pandeng}

\subsection{机器当望远镜：探索几何、数论与分析}

计算机最浪漫的角色，是做数学家的"望远镜"和"显微镜"——帮人类\textbf{看见}原本看不见的数学世界，至于证明，仍然由人来完成。

在\textbf{几何}里，计算机把迭代过程画成图，炸出了一个全新的世界。曼德博集合——按规则$z \mapsto z^2 + c$对复数反复迭代，把"永远不跑远"的点涂黑——画出来的图形边缘层层分形，放大任何一小段边缘都会涌现出新的花纹，无穷无尽。在计算机出现之前，没有人知道如此简单的公式里藏着如此繁复的几何。今天"分形几何"已经用来描述海岸线、云和血管的形状。

在\textbf{数论}里，计算机是超级筛子。已知的最大素数是梅森素数，有两千多万位，由全球志愿者用闲置电脑联网搜寻（GIMPS 计划）找到。人们还用计算机验证哥德巴赫猜想（"每个大于 $2$ 的偶数都是两个素数之和"）到极大的范围、统计素数分布的精细规律——这些海量数据常常为猜想提供灵感，有时也为推翻猜想提供弹药。

在\textbf{分析}里，数值实验时常先于理论发现规律。一个传奇的例子是 $\pi$ 的 BBP 公式：人们用计算机做数值搜索，先发现了一个能从 $\pi$ 的十六进制展开中直接"抠出"任意一位的公式，随后才补上严格证明。机器先"看"到规律，人再搞懂它为什么成立——探索的顺序被颠倒了，而这正是实验科学的经典流程。

\begin{toolbox}
\textbf{本章数学工具箱}

\begin{itemize}
\item \textbf{有限验证 $\neq$ 证明}：验证再多有限情形也只提供证据；一个反例足以推翻猜想，一个证明必须覆盖无穷情形。
\item \textbf{算法}：有限、明确、必然终止的指令序列。它是讨论"能不能机械求解"的精确语言。
\item \textbf{停机问题}：存在\textbf{原则上}无法用算法判定的问题；自我指涉 + 反证法是识别这类障碍的利器。
\item \textbf{计算复杂度}：在有算法的问题中再分快慢；多项式增长可控，指数增长是怪兽；P 对 NP 是最大的未解之谜之一。
\item \textbf{人机分工}：机器做海量计算、严格校对、图像探索；人负责提出问题、发明概念、完成真正需要洞察的推理。
\end{itemize}
\end{toolbox}

\section{纸上机房：动手实验}

\begin{shiyan}
\textbf{实验一：亲手跑一遍 $3n+1$。}

从 $27$ 出发，在纸上一步一步执行规则：偶数除以 $2$，奇数乘 $3$ 加 $1$。耐心算下去，数一数用了多少步才到达 $1$（答案在一百步上下，做好心理准备），并记录途中出现的最大值。你会发现 $27$ 这个小个子居然能蹿到 $9232$ 那么高再跌下来。再试试从 $31$ 或 $97$ 出发。体会一下：为什么计算机验证这个猜想并不轻松——路径忽高忽低，毫无规律，却\textbf{至今}每次都有惊无险地落地。

\textbf{实验二：用一摞书感受二分查找。}

把 $16$ 本书在桌上按书脊的高矮（或页码）从左到右排好。请一位同学心里默想其中一本，你只许问"是这本吗，还是说目标在它的左边/右边"。先每次从中间那本问起，记录你最多几次就能锁定目标；再换从第一本开始挨个问，记录最坏情况要几次。$16$ 本书，二分法最多 $4$ 次，挨个问最坏 $16$ 次。把书加到 $32$ 本重做一次：二分法只多 $1$ 次，挨个问却翻了一倍——这就是对数增长和线性增长的差别，亲手摸一遍就再也忘不掉。

\textbf{实验三：给"撒谎的程序"画像。}

不碰电脑，只在纸上推演：假设同桌宣称他写了一个万能程序 $H$，能判断任何程序会不会死循环。模仿本章定理的证明，向同桌描述那个专门跟他对着干的程序 $D$，让他亲口说出"$D$ 作用在 $D$ 上到底停不停"。无论他怎么回答，你都指出其中的矛盾。这个纸笔游戏做完，你会对"反证法"和"自我指涉"产生真正的手感。
\end{shiyan}

\section*{思考题}

\textbf{观察题}

1. 用 $n^2 - n + 41$ 对 $n=1,2,\dots,8$ 算出结果并检查是否都是素数，再直接计算 $n=41$ 的情形。（提示：$n=41$ 时前两项 $41^2-41$ 能提取出公因数 $41$，所以整个式子必有因数 $41$。）

2. 观察 $3n+1$ 过程中：任何一个到达 $1$ 的轨道，最后几个数必定是什么？（提示：想想什么样的数经过"除以 $2$"能变成 $1$，再往前倒推一步。）

\textbf{证明题}

3. 模仿停机问题的证明，说明不存在这样的算法 $G$：给它任意程序 $P$，$G$ 总能正确判断"$P$ 运行时会不会在屏幕上打印出'你好'两个字"。（提示：假设 $G$ 存在，造一个怪程序：先问 $G$"我会打印你好吗"，听到"会"就偏不打印，听到"不会"就偏要打印，然后让它判断自己。）

4. 证明：如果书架上有 $n$ 本按顺序排好的书，二分查找最多需要 $\lceil \log_2 n \rceil + 1$ 次以内就能找到任意一本。（提示：每翻一次，待查范围至少减半；想想范围从 $n$ 缩到 $1$ 需要减半多少次。）

\textbf{挑战题}

5. 有人宣布他证明了"P 等于 NP"，并且找到了在多项式时间内解决数独的算法。假设他是对的，列举三件我们日常生活中会因此天翻地覆的事，并说明理由。（提示：回忆本章讲的密码学；再想得更远一点——数学定理的"证明"本身，验证起来难不难？如果验证容易的事寻找也容易，"自动证明定理"意味着什么？）

6. $3n+1$ 猜想至今没被证明。有人说："既然停机问题不可判定，说不定 $3n+1$ 猜想本质上也是'不可判定的'，所以才证不出来。"这个推理有哪些地方站得住，哪些地方是偷换概念？（提示：注意"不存在一个万能算法判定\textbf{所有}程序"和"某\textbf{一个}具体命题无法被证明"是两回事；但也确实有数学家构造出与停机问题本质等价的具体数论问题。）

\section*{通往下一章的门}

这一章我们看到了一幅奇特的图景：计算机是人类造过的最强大的计算工具，它能在眨眼间验证三亿亿个数，能逐行检查人一辈子查不完的证明，能画出人脑无法想象的分形——可是它面前同样矗立着不可逾越的高墙：有些问题永远没有算法，有些问题算法慢到等不起，而"验证了一亿次"在数学上依然不等于"永远成立"。

机器的力量和机器的边界，其实都映照出数学自身的力量和边界。哥德尔说不完备性是任何足够强的公理系统逃不掉的宿命；图灵说不可判定性是任何计算装置绕不开的暗礁；而"P 对 NP"这样的谜题至今悬而未决，提醒我们今天引以为傲的大厦依然建在不少未探明的地基上。

那么，数学这座大厦的地基究竟是什么样子？极限、无穷、连续这些让初中数学开始"不踏实"的概念，是怎样被十九世纪的数学家们一寸一寸夯实的？那些深刻的定理——从微积分基本定理到伽罗瓦理论——又是怎样把看似毫不相干的数学分支缝合成一张统一地图的？下一章，我们将走到数学世界的边缘，回头看它的全貌：边界在哪里，地图又通向何方。
